% Section 1 -- Introduction. Budget: 1.00 page, holds Figure 1.
\section{Introduction}
\label{sec:intro}

Written Vietnamese is a syllabic script in a Latin alphabet: every syllable is written as a separate word, and its three parts are visible on the page. The onset is a consonant letter or digraph, the rime is a vowel sequence with an optional final consonant, and the tone is one of six values marked by a diacritic (or by its absence). A literate speaker can take a syllable apart at these seams and put it back together, and Vietnamese has a folk art that does exactly that. \vi{Nói lái} swaps the rimes, the tones, or both between two syllables of a phrase: \vi{mèo cái} `female cat' becomes \vi{mài kéo} `sharpen scissors', and \vi{bí mật} `secret' becomes \vi{bị mất} `got lost'. The transformation is defined entirely below the level of the word, and the spelling of the result follows rules that depend on the new neighbours: when the rime \vi{eo} moves next to the onset /k/, the letter changes from \vi{c} to \vi{k}.

Large language models do not see any of this structure. Subword tokenizers split Vietnamese syllables where their merge statistics tell them to, not at the onset--rime boundary, and when text arrives in a different but equally valid Unicode encoding the split changes altogether. Figure~\ref{fig:tokens} shows the Gemma~3 tokenizer on the example above: in the precomposed (NFC) encoding the output syllable \vi{mài} is cut between onset and rime, and in the combining (NFD) encoding the tone mark of \vi{mèo} becomes a token of its own. Our audit of the same tokenizer over the 6,595 standard Vietnamese syllables finds that it splits 70\% of them under NFC, half of all syllables exactly at the onset--rime boundary, and spends one extra token per syllable under NFD, where the tone mark is isolated in 42\% of syllables.\footnote{Numbers from the tokenizer audit in our released code (\texttt{data/audit/gemma3.json}); the tokenizer does not normalize NFD to NFC, so the two encodings reach the model as different token sequences.} The question this paper asks is therefore general: \emph{can a subword LLM reach the onset, rime and tone inside a tonal syllable it never receives as a unit, and does failure grow with how badly the tokens cut across those units?}

Nói lái is our instrument, not our headline. Because the documented variants are permutations of syllable components and Vietnamese spelling is rule-governed, a small rule engine can generate transformation, decoding and validity items in unlimited quantity, verify every gold answer, and score every model output by parsing it back into onset, rime and tone. No LLM judge is needed anywhere. The same engine lets us re-encode any Vietnamese text between canonically equivalent Unicode forms and between the two tone-mark placement conventions while holding meaning and information fixed, which turns deployed tokenizers into the object of a training-free intervention. Finally, the two smallest Gemma~3 models let us ask where in the network tone information is represented, whether it reaches the answer position, and whether it is used, following the template of character-level probing work \citep{2025-findings-emnlp-719}.

% The caption's data sentences (\figtokensnote) are GENERATED by gen_fig1_tokens.py from the
% same audit rows as the table, so that the caption cannot contradict the numbers it sits under.
% GENERATED by paper/gen_fig1_tokens.py from /home/user/noilai-migrate/data/external/gemma3_tokenizer.model on 2026-10-01 13:09Z; do not edit by hand.
% \figtokensnote: the caption sentences of Figure 1 that state what the token table shows,
% composed from the audit rows so that the caption cannot contradict the figure.
\newcommand{\figtokensnote}{Under NFC \vi{mèo}, \vi{cái} and \vi{kéo} are single tokens and \vi{mài} is cut at the onset$|$rime seam (alignment 1.00). Under NFD every syllable gains a boundary between a base letter and its combining tone mark, which is never linguistic, so alignment falls to 0.50 for \vi{mèo} and \vi{kéo}, whose tone mark becomes a token of its own, leaving one aligned boundary at the nucleus$|$coda seam, and to 0.00 for \vi{cái} and \vi{mài}, whose tone mark stays attached to the following letter.}

\begin{figure}[t]
  \centering
  \small
  % GENERATED by paper/gen_fig1_tokens.py from /home/user/ntcf/noilai/data/external/gemma3_tokenizer.model on 2026-09-30 15:10Z; do not edit by hand.
% tokenizer: BPE 262,144 pieces, normalizer identity, byte_fallback True; syllables tokenized in running-text position (leading space).
\begin{tabular}{@{}lllcll@{}}
\toprule
 & \vi{mèo} & \vi{cái} &  & \vi{mài} & \vi{kéo} \\
 & \multicolumn{2}{l}{input (\vi{mèo cái})} & $\xrightarrow{\ \var{1}\ }$ & \multicolumn{2}{l}{output (\vi{mài kéo})} \\
\midrule
NFC (precomposed) & \tok{\tokspace{}mèo} & \tok{\tokspace{}cái} &  & \tok{\tokspace{}m}\tok{ài} & \tok{\tokspace{}kéo} \\
\quad tokens / alignment & 1 / 1.00 & 1 / 1.00 &  & 2 / 1.00 & 1 / 1.00 \\
NFD (combining) & \tok{\tokspace{}me}\tok{\`{}}\tok{o} & \tok{\tokspace{}ca}\tok{\'{}i} &  & \tok{\tokspace{}ma}\tok{\`{}i} & \tok{\tokspace{}ke}\tok{\'{}}\tok{o} \\
\quad tokens / alignment & 3 / 0.50 & 2 / 0.00 &  & 2 / 0.00 & 3 / 0.50 \\
\bottomrule
\end{tabular}

  \caption{\vi{Mèo cái} $\rightarrow$ \vi{mài kéo} (variant \var{1}: the rimes swap, onsets and tones stay) with the token boundaries of the Gemma~3 tokenizer, each syllable in running-text position. \protect\tokspace{} marks the SentencePiece word-initial space; an accent over an empty box is a combining mark that forms a token on its own. Alignment is the share of internal token boundaries that fall on an onset$|$rime, glide$|$nucleus or nucleus$|$coda seam; a boundary between a base letter and its combining mark never counts. \figtokensnote{} \placeholder{Add model outputs for three tokenizers once E1 has run.}}
  \label{fig:tokens}
\end{figure}

We make five contributions.
\begin{enumerate}[nosep,leftmargin=*]
  \item \textbf{\noilai{}}, the first LLM study of nói lái and the first generative, decoding and validity benchmark of it: a rule-verified generator over the four generated variants of a six-type tradition, an attested cultural subset, native-speaker validation and a human baseline (Section~\ref{sec:benchmark}).
  \item \textbf{A syllable-alignment audit} that links how production tokenizers split onset, rime and tone to item-level accuracy across \placeholder{21} models (Sections~\ref{sec:setup} and~\ref{sec:results}).
  \item \textbf{Orthographic counterfactuals} that change only the token sequence of a tonal Latin script on deployed models, separating tokenizer-caused from knowledge-caused errors, with a census of which tokenizers normalize, pass through or corrupt combining marks (Section~\ref{sec:counterfactuals}).
  \item \textbf{A mechanistic localization} of tone information in small open models, using control-validated probes, structural baselines and activation patching with separate perception and manipulation readouts (Section~\ref{sec:tone}).
  \item \textbf{Reusable tooling}: the syllable parser, spelling-rule engine and re-encoder apply to any diacritic-heavy Latin script.
\end{enumerate}

Every headline number carries a 95\% interval that resamples base pairs, every comparison is a paired test corrected within its family, every claim about models in general rests on a pooled estimate rather than a count of significant cells, and the hypotheses \hyp{1}--\hyp{6} named below were pre-registered in a dated commit before any test-split run (the pilot runs on development items only). \placeholder{One-sentence summary of the main finding, written from the results: TODO.}
